
\documentclass[final,ignorenonframetext,compress]{beamer}

\mode<beamer>
{
\usetheme{Bonn}
\pgfdeclareimage[height=0.8cm]{university-logo}{images/Logo_UBo_h24_4c-crop}
\setbeamertemplate{footline}[default]
\logo{\pgfuseimage{university-logo}}
\setbeamertemplate{navigation symbols}{}
}
\mode<handout>{
\usetheme{default}
\usecolortheme{dove}
\setbeamercolor{background canvas}{bg=black!5}
\usepackage{pgfpages}
\pgfpagesuselayout{8 on 1}[a4paper,border shrink=5mm]
}
\usepackage{natbib}
\usepackage{booktabs}
\usepackage{tabularx}

\DeclareMathOperator*{\mymax}{max}
\DeclareMathOperator*{\argmax}{arg\,max}
\DeclareMathOperator*{\argmin}{arg\,min}
\DeclareMathOperator*{\ForAll}{\bigwedge}
\DeclareMathOperator*{\Exists}{\bigvee}
\newcommand{\W}{\mathcal{W}}

\newcommand*\kreis[2]{\unitlength0.7ex\begin{picture}(2.5,2.5)%


	\color{#2}%
	\put(0.75,0.75){\circle*{3}}%
	\color{black}%
	\put(0.75,0.75){\circle{3}}%
	\put(0.7,0.7){\makebox(0,0){#1}}\end{picture}}
	\newcommand\itplus{\kreis{$+$}{green}}
	\newcommand\itminus{\kreis{$-$}{red}}
	\newcommand\prob[2]{\ensuremath{p_{\mathrm{#1}}\left( #2 \right)}}
	\newcommand\bedprob[3]{\ensuremath{p_{\mathrm{#1}}\left( #2 \left|#3\right. \right)}}

    \newcommand\etal{et.\,al.}

	\usepackage[english]{babel}
	% avant, courier, chancery, times, palatino, bookman, 
	% newcent, utopia, charter
	%\usepackage{times}

	\usepackage[latin1]{inputenc}
	\usepackage[T1]{fontenc}
	\usepackage{multicol}
	\usepackage{mdwtab}
	\usepackage{booktabs}
    \usepackage{tabularx}
	\usepackage{txfonts,textcomp}
	\usepackage{amsmath}
	\usepackage{amsfonts}

	\setlength\parindent{0cm}
	\setlength\parskip{2mm}
	\setbeamerfont{gross}{size=\large}
	\setbeamerfont{quetsch}{size=\tiny}
	\setbeamerfont{klein}{size=\footnotesize}
	\newcommand\w[1]{\mathbf{#1}}

	% example list spacing
	\newcommand\ml{\setlength\itemindent{-11mm}}

    \newcommand{\newblock}{}


	\title{
Learning Depth-Sensitive Conditional Random Fields for Semantic Segmentation of RGB-D Images
    } 
    \author[short]{Andreas M\"uller}%
	\institute[UBO]{%
    Advisor: Sven Behnke\\
    Institute for Autonomous Intelligent Systems\\
	Rheinische Friedrich-Wilhelms-Universit\"at\\[4mm]
	\includegraphics[width=.2\linewidth]{images/Logo_UBo_h24_4c-crop}\\[3mm]
	Bonn, Germany}


	\begin{document}

	\begin{frame}[plain]
		\titlepage
	\end{frame}

    \section*{Introduction}

	\begin{frame}
		\frametitle{Semantic Segmentation}
            \begin{figure}
                \includegraphics[width=.4 \textwidth]{images/pascal}
                \begin{visibleenv}<2->
                    \includegraphics[width=.4 \textwidth]{images/pascal_gt}
                \end{visibleenv}
            \end{figure}
        \begin{itemize}
            \item<3-> Label each pixel with semantic class.
            \item<4-> Important step towards scene understanding.
        \end{itemize}

	\end{frame}

    %\begin{frame}
        %\frametitle{Semantic Segmentation using Super Pixels}
        %\begin{figure}
            %\includegraphics<1>[width=.6\linewidth]{images/scene_sp_org}
            %\includegraphics<2->[width=.6\linewidth]{images/scene_sp}
        %\end{figure}
        %\vspace{-5mm}
        %Grouping pixels into superpixels:
        %\begin{itemize}
            %\item Reduces the computational burden.
            %\item Provides more context for classification.
        %\end{itemize}
    %\end{frame}

    \section{Structured prediction formulation}

    \begin{frame}
        \frametitle{Semantic Segmentation as Structured Prediction}
        \begin{figure}
            \includegraphics<1-2>[width=.6\linewidth]{images/scene_sp_gt}
            \includegraphics<3->[width=.6\linewidth]{images/scene_sp_graph}
        \end{figure}
        \vspace{-5mm}
        \begin{itemize}
            \item<1->Predict class for each superpixel.
            \item<2->Predictions clearly not independent.
            \item<3->Use neighborhood structure to model correlations.
        \end{itemize}
    \end{frame}

    \begin{frame}
        \frametitle{Structured Prediction}
        Learn prediction function of the form
        \[f(x, w) := \argmax_{y \in \mathcal{Y}}  w^T \psi(x, y) \]
        \begin{visibleenv}<2->
        where
        \begin{itemize}
            \item $x$ is the input.
            \item $w$ are the parameter of the model to be learned.
            \item $\psi(x, y)$ is a joint feature function of $x$ and potential labeling $y$.
            \item $\mathcal{Y}$ is the output space (often exponentially large).
        \end{itemize}
        \end{visibleenv}
    \end{frame}

    \section{Model and Features}
    \begin{frame}
        \frametitle{Superpixels}
        % image + formular!
    \end{frame}

    \begin{frame}
        \frametitle{Pairwise Models}
        Typically $\psi$ decomposes over pairs of entries $(y_i, y_j)$:
        \[w^T \psi(x, y) = \sum_{(i, j) \in E} w_{i,j} \psi_{i,j}(x, y_i, y_j) + \sum_i w_i \psi_i(x, y_i)\]
        \vspace{-8mm}
        \begin{visibleenv}<2->
            \begin{center}
                \includegraphics[width=.5 \linewidth]{images/scene_sp_graph}
            \end{center}
        \end{visibleenv}
        \vspace{-5mm}
        \begin{visibleenv}<3->
            In loopy models, exact inference is intractable in general, but many approximations are available.
        \end{visibleenv}
    \end{frame}

    \begin{frame}
        \frametitle{Random Forest Unary Potentials}
    \end{frame}

    \begin{frame}
        \frametitle{Absolute height potentials}
        \begin{center}
            \includegraphics[width=\linewidth]{images/height_success}\\
            \vspace{3mm}
            \includegraphics[width=\linewidth]{images/height_failure}
        \end{center}
    \end{frame}


    \begin{frame}
        \frametitle{Color Contrast}
    \end{frame}

    \begin{frame}
        \frametitle{Vertical Alignment}
    \end{frame}

    \begin{frame}
        \frametitle{Depth Difference}
    \end{frame}
    

    \begin{frame}
        \frametitle{Pairwise Normal Orientation}
        \includegraphics[width=\linewidth]{images/normal_feature}
    \end{frame}


    \begin{frame}
        \frametitle{Summary of Pairwise Potentials}
            \noindent\emph{$\bullet$~Constant.}
                A constant feature allows to model general neighborhood relations.


            \noindent\emph{$\bullet$~Color Contrast.}
                We employ a non-linear color contrast, as is common in the computer
                vision literature, between the superpixel mean colors $c_i$ and $c_j$:
                $\exp\left(-\gamma \lVert c_i - c_j \rVert^2\right)$.


            \noindent\emph{$\bullet$~Vertical Alignment.}
                We model the directed angle between superpixel centers in the 2D image
                plane.  This allows the  model to learn that ``structure'' is above
                ``floor'', but not the other way around.


            \noindent\emph{$\bullet$~Depth Difference.}
                We include the signed depth difference between superpixels, which
                allows the model to detect depth discontinuities that are not
                represented in the 2D neighborhood graph of the superpixels.


            \noindent\emph{$\bullet$~Normal Orientations.}
                Differences in normal vector orientation are a strong clue on
                whether two superpixels belong to the same surface, and therefore the
                same structural class.
                We compute the 3D orientation of normals using the method of \citet{holz_2011_robocup},
                as implemented in the point cloud library (pcl)\footnote{\url{http://pointclouds.org}}.
                All normals within a superpixel are then averaged, to get a single orientation for each superpixel.
                The feature is computed as the difference of $\frac{\pi}{4}$ and the (undirected) angle between the normals belonging
                to two adjacent superpixels.  
                An example is shown in \Figref{normal_feature}. The change
                in normal orientation highlights that pillow and wall are distinct
                objects, even though there is no strong distinction in color or depth.
    \end{frame}

    %\begin{frame}
        %\frametitle{Solving the Structured SVM problem}
        %\emph{Inference} is critical for learning:

        %\begin{align*}
            %&\argmax_{y \in Y} \left < w, \psi(x_i, y) \right > &\text{(prediction)}\\
            %&\argmax_{y \in Y} \left < w, \psi(x_i, y) \right > + \Delta(x_i, y_i, y) &\text{(loss-augmented prediction)}
        %\end{align*}
        %\begin{visibleenv}<2->
             %Typically $Y$ is (exponentially) large, making enumeration impossible.
         %\end{visibleenv}

        %\begin{visibleenv}<3->
             %We need to exploit the \emph{structure} of the problem, encoded in $\psi$.
         %\end{visibleenv}
    %\end{frame}

    %\begin{frame}
        %\frametitle{Pairwise models as Structured SVMs}
        %\begin{columns}
            %\column{.6\linewidth}
            %\begin{itemize}
                %\item<1-> Typically $\psi$ decomposes over pairs of entries $(y_i, y_j)$:
                    %\[w^T \psi(x, y) = \sum_{(i, j) \in E} w_{i,j} \psi_{i,j}(x, y_i, y_j) + \sum_i w_i \psi_i(x, y_i)\]
                %\item<2-> If $E$ is a tree, exact inference by message-passing is possible.
                %\item<3-> If $E$ does contain circles, we can still do approximate
                    %inference using
                    %\begin{itemize}
                        %\item message passing
                        %\item graph cuts
                        %\item linear programming 
                        %\item dual decomposition
                    %\end{itemize}
            %\end{itemize}
            %\column{.3\linewidth}
            %\vspace{40mm}
            %\includegraphics[width=\linewidth]{images/scene_sp_graph}
        %\end{columns}
    %\end{frame}

    %\begin{frame}
        %\frametitle{Solving the Structured SVM problem}
        %Once we know how to do inference, there are several approaches to finding $w$:
        %\begin{itemize}
            %\item Dualizing the problem, leading to a quadratic program \citep{taskar2003max, tsochantaridis2006large, joachims2009cutting}
            %\item Stochastic Subgradient Descent \citep{ratliff2007online, lucchi2013learning}.
            %\item Online coordinate descent in the dual \citep{lacoste2012block, shalev2012proximal, branson2013efficient}.
            %\item Combining inference and learning \citep{meshi2010learning, hazan2010primal, stoyanov2011empirical, jancsarylearning}.
        %\end{itemize}
    %\end{frame}

    \begin{frame}
        \frametitle{How intractable are loopy models?}
        \begin{itemize}
            \item In general, inference in loopy models is NP-hard.
            \item Learning relies on inference.
            \item Combining several inference algorithms, we can get exact results in practice. %FIXME reference
        \end{itemize}
    \end{frame}


    %\begin{frame}
        %\frametitle{Conditional Random Fields and Structured SVMs}
        %\begin{itemize}
            %\item Conditional Random Fields: model $p(y | x)$, predict most likely $y$.
            %\item Structured SVM: model $f(x, y)$, predict $\argmax_y f(x, y)$.
            %\item Optimizing $p(y| x)$ is intractable, optimizing $L(\argmax_y f(x, y), y)$
                %slightly less so.
            %\item Computer vision researchers still call all graph based structured models CRFs.
        %\end{itemize}
    %\end{frame}


    %\begin{frame}
        %\frametitle{Applying Structured SVMs to Image Segmentation}
        %\begin{itemize}
            %\item<1-> Data-dependend terms $\psi_i(x, y_i) = x_i \otimes e_{y_i}$.
            %\item<2-> Often $\psi_{i, j}(x, y^i, y^j) = e_{y_i} \otimes e_{y_j}$.
            %\item<3-> \emph{Strength} of interaction sometimes conditioned on image:
                %\[\psi_{i, j}(x, y^i, y^j) =  \alpha(||x_i - x_j||) \cdot e_{y_i} \otimes e_{y_j}.\]
        %\end{itemize}
    %\end{frame}

    \section{Experiments}
    
    \begin{frame}
        \frametitle{Dataset}
    \end{frame}
    

    \begin{frame}
        \frametitle{Experimental Protocol}
    \end{frame}

    \begin{frame}
        \frametitle{Results}
        \begin{table*}[t]
            \caption{Quantitative comparison of the proposed method with the
            literature.\tablabel{results}}
        \begin{tabularx}{\linewidth}{@{\extracolsep{\fill}}lcccccc}
        \toprule
                                & ground        &  structure    & furniture     & props         & class average   & pixel average\\
        \cmidrule(r){1-7}

        RF                              &         90.8  &   81.6        & 67.9          & 19.9          &  65.0        &  68.3 \\
        RF + SP                         &         92.5  &   83.3        & \textbf{73.8} & 13.9          &  65.7        &  70.1 \\ 
        RF + SP + SVM                   &         94.4  &   79.1        & 64.2          & \textbf{44.0} &  70.4        &  70.3 \\
        RF + SP + CRF                   & \textbf{94.9} &   78.9        &          71.1 & 42.7          &\textbf{71.9} &  \textbf{72.3} \\
        \cmidrule(r){1-7}
        \citet{SilbermanECCV12}         &         68    &   59          & 70           & 42            &  59.6        & 58.6 \\
        \citet{couprie-iclr-13}         &         87.3  & \textbf{86.1} & 45.3         & 35.5          &  63.5        & 64.5 \\
        \cmidrule(r){1-7}
        \citet{stueckler2013}$^\dagger$ & \textbf{95.6} &   83.0        & \textbf{75.1}& 14.2          &  67.0        & 70.9 \\

        \bottomrule
        \end{tabularx}
        \quad\\The best value in each column is printed in bold$^\dagger$. The upper part of
        the table shows contributions by different parts of our pipeline. RF stands for random forest prediction, RF + SP for aggregated
        random forests prediction within superpixels, RF + SP + SVM for an SVM trained on the unary potentials, and RF + SP + CRF is
        our proposed pipeline. We optimized our approach for class average
        accuracy.\\
        $^\dagger$ \footnotesize Note that the work of \citet{stueckler2013} is not directly
        comparable, as they integrated information over multiple frames, and did not
        measure accuracy for pixels without valid depth measurement.
        \vspace{-5mm}
        \end{table*}
    \end{frame}

    \begin{frame}
        % qualitative results
    \end{frame}


    \section*{Structured Machine Learning in Python}
    \begin{frame}
        \frametitle{PyStruct - simple structured prediction}
        \begin{columns}
            \column{.6\textwidth}
                \begin{itemize}
                    \item Implements common solvers: perceptron, cutting plane
                        ssvm, latent ssvm, subgradient, block-coordinate
                        Frank-Wolfe.
                    \item<2-> Implements powerful models for many use cases.
                    \item<3-> Interface to popular inference libraries: LibDAI, OpenGM, QPBO, AD3.
                    \item<4-> Lots of documentation and examples.
                \end{itemize}
            \column{.4\textwidth}
            \begin{center}
                \includegraphics[width=.6\linewidth]{images/logo_pystruct}
            \end{center}
        \end{columns}
    \end{frame}

    \begin{frame}
        \begin{center}
            \includegraphics[width=\linewidth]{images/example_gallery}\\
            http://pystruct.github.io
        \end{center}
    \end{frame}

    \begin{frame}
        \frametitle{Take Home}
        \begin{itemize}
            \item High level semantic unaries.
            \item Learn everything.
            \item Exact learning of CRF possible.
            \item Straight forward extension to 3D features.
        \end{itemize}
    \end{frame}
    
    \section*{Thank you.}
    \begin{frame}
        \begin{center}
            Thank you for your attention.

            Questions?
        \end{center}
    \end{frame}
    \begin{frame}<beamer:0>
        \bibliographystyle{plainnat}
        \bibliography{paper}
    \end{frame}

	\end{document}
